\originalsection{完整标准 Transformer 的 CPU 配对训练}\label{sec:standard-model}
手算展示一条注意力路径，完整模型还要把位置、归一化、前馈、残差和读出连接起来。\path{standard_transformer_study.py} 使用 PyTorch 的标准层实现两层前置归一化的因果解码器；它不是包装外部 API，也没有下载预训练权重。结构如下：词表 10、宽度 $32$、两头、前馈宽度 $64$、上下文长度 7，两个残差块，GELU 激活，最后归一化及线性读出，总计 18,026 个可训练参数。位置是可学习嵌入，没有使用本章另一节推导的 RoPE；实现和方法名不能混淆。

每一块可写成
\[
U=X+\operatorname{MHA}(\operatorname{LN}(X)),\qquad
X'=U+\operatorname{FFN}(\operatorname{LN}(U)),
\]
其中 $\operatorname{FFN}(u)=\operatorname{GELU}(uW_1+b_1)W_2+b_2$。训练对整个词表用交叉熵，不只监督叫声。八条原创语法句从说话者、颜色、旁观者各两种组合构成；固定六条训练、两条留出，整个设计在训练前写成独立 JSON 并计算哈希。

我们比较均衡六条训练句与将同一条训练句额外重复 30 次。若单条损失是 $\ell_i$，第二组目标为
\[
L_{\rm repeat}=\frac{31\ell_0+\sum_{i\ne0}\ell_i}{36},
\qquad L_{\rm balanced}=\frac16\sum_{i=0}^5\ell_i.
\]
唯一句子数没有增加，分布权重改变了。两组同一种子从相同参数出发、做 600 次 AdamW 更新，但每步分别处理 6 与 36 行，因此没有匹配 FLOPs。这个实验检验频率重加权及其伴随处理成本，不能当成等算力数据质量消融，更不是 FineWeb 的大规模复现。

种子固定为 144、145、146，学习率 $0.003$、权重衰减 $0.01$、梯度裁剪 1，无 dropout，CPU 单线程 float32。六次运行全部报告，未用留出成绩挑最好种子。下面由实际结果 JSON 自动生成，代码与环境哈希另附。
\begin{center}\begin{tabular}{clrrr}\toprule
种子 & 条件 & 加权训练 NLL & 均衡训练 NLL & 留出 NLL \\ \midrule
144 & 均衡 & 0.2565 & 0.2565 & 1.3159 \\
144 & 单条重复 & 0.0899 & 0.4310 & 1.3747 \\
145 & 均衡 & 0.2565 & 0.2565 & 1.2958 \\
145 & 单条重复 & 0.0900 & 0.4311 & 1.3475 \\
146 & 均衡 & 0.2565 & 0.2565 & 1.2773 \\
146 & 单条重复 & 0.0900 & 0.4311 & 1.3985 \\
\bottomrule\end{tabular}\end{center}

重复条件的自身加权训练损失更低，三个配对留出损失却更高。若只看优化器监控的训练目标，会错把偏向常见句的改进当作总体泛化提高；使用同一均衡训练分布评价又得到不同结果。这里两条留出样本太少，不能从三次初始化推出广泛的数据结论。

六个模型对原始、改变旁观者及改变说话者三种叫声测试均符合该有限语法；但这些干预组合与训练句子重叠，不能识别唯一内部机制，也不能证明任意组合迁移。编辑未来 token 后先前位置的最大 logits 差为零；额外测试还检查了未来嵌入的梯度为零。这提供了因果掩码实现证据。想继续研究，应事先扩充组合、独立留出角色和词表，匹配 token 处理预算，再比较数据重加权；不能看到当前结果后悄悄修改这一次已冻结实验。
